% Empirical strategy. States no result. Source: plan.json (hypotheses, design.estimand,
% identifying_assumption, method, analysis_plan, threats, deviations). Every number below is a
% fact of the approved design. The check-methods outcomes that shaped the design (design.method
% C1-C3, T6) are described qualitatively here; their values enter the paper only as macros of
% claims backed by the run level-mde, which copies them (design.analysis_plan 2, claims_plan).
\section{Empirical strategy}
\label{sec:strategy}

The design compares California's limited-service restaurant employment after AB 1228 took effect
with a counterfactual built from other states. Which counterfactual---which estimator, donor pool,
weighting and outcome normalisation---is not chosen by judgement: a rule fixed at the design gate
picks it from 24 candidates using only data that end in 2019Q4. Inference is built for a single
treated state, and the labels the paper reports follow from decision rules fixed at the same time.

\subsection{Estimand}
\label{sec:estimand}

Let $Y_{\mathrm{CA},t}$ be the seasonally adjusted log employment of private NAICS 722513
establishments in California in quarter $t$, and $Y_{\mathrm{CA},t}(1)$ and $Y_{\mathrm{CA},t}(0)$
its values with and without AB 1228. The estimand is
\begin{equation}
\label{eq:estimand}
\tau = \frac{1}{8}\sum_{t=\text{2024Q2}}^{\text{2026Q1}}
  \bigl[Y_{\mathrm{CA},t}(1) - Y_{\mathrm{CA},t}(0)\bigr]
  - \bigl[Y_{\mathrm{CA},\text{2023Q3}}(1) - Y_{\mathrm{CA},\text{2023Q3}}(0)\bigr].
\end{equation}
If AB 1228 had no effect before it was signed (September 28, 2023, after the September pay period
that the QCEW counts), the second bracket is zero. $\tau$ is an intent-to-treat effect on the whole
industry code, covered and uncovered establishments together (Section~\ref{sec:coverage}), in log
points; we also report it as $100(e^{\tau}-1)$ percent and as jobs, the mean of California's
published (not seasonally adjusted) 2024Q2--2026Q1 employment times $1 - e^{-\tau}$. The wage
estimand of hypothesis H1 is the same with the log average weekly wage. Hypothesis H3 concerns
$\tau_2 - \tau_1$, where $\tau_1$ averages over year one (2024Q2--2025Q1) and $\tau_2$ over year two
(2025Q2--2026Q1). Candidates fitted on log employment per capita estimate the same $\tau$ under the
maintained assumption that AB 1228 did not change California's resident population, which is what
lets the selection rule compare all candidates on one scale.

\subsection{Timing}
\label{sec:timing}

The base quarter is $B = \text{2023Q3}$, the last quarter before AB 1228 was signed. The two quarters
between signing and the effective date, 2023Q4 and 2024Q1, form an anticipation window that enters
neither the base nor the post window; we report the effect in that window separately, without a
test, as \citet[Table~A3]{sosinskiy2026effects} already do at county level. The COVID quarters
2020Q1--2021Q4 are excluded from every fit, base, trend and standard error, as in
\citet{dube2026labor}, with the one exception noted for the secondary analysis S1
(Section~\ref{sec:secondary}). In general, for a base quarter $B$,
\begin{itemize}
  \item the fit window $F(B)$ is every quarter from 2014Q1 to $B$ outside the COVID quarters;
  \item the anticipation window is $A(B) = \{B+1, B+2\}$;
  \item the post window $P(B)$ is the eight quarters $B+3, \dots, B+10$, with year one $B+3, \dots, B+6$
        and year two $B+7, \dots, B+10$.
\end{itemize}
For the real base, $F$ is 2014Q1--2019Q4 and 2022Q1--2023Q3 (31 quarters), $A$ is 2023Q4--2024Q1 and
$P$ is 2024Q2--2026Q1. Time is the calendar quarter index, so the COVID gap counts as elapsed time
in trends. The published result closest to this timing is \citet[Table~1, Panel~B]{dube2026labor},
whose base, the 2022Q1--2023Q3 average, also keeps 2023Q4 out. Dube's own 2023Q4 base and a base
before AB 257 (2022Q2) are pre-specified robustness checks (Section~\ref{sec:robust-plan}), not
alternative primaries.

\subsection{Identifying assumption}
\label{sec:assumption}

Without AB 1228, California's seasonally adjusted log 722513 employment (for H1, its log weekly
wage) after 2023Q3 would have stayed as far from the selected counterfactual as it was before. For
the DD candidates the reference is the 2023Q3 gap, with the pre-period trend of the gap extrapolated
for the trend-adjusted variant; for synthetic control it is the mean gap over the fit window; for
SDID it is the time-weighted pre-period gap.

The selection rule rests on a premise: the candidate that best reproduced California's actual path
at pre-COVID pseudo-treatment dates is the one for which the assumption holds best. That path
includes California's own general minimum-wage ramp (Section~\ref{sec:ramp}), which is part of the
counterfactual for AB 1228. The premise is therefore about fit to California's path under its own
ramp, and it assumes that whatever lets a candidate track that ramp in 2016--2019---a gap trend, or
donors with their own increases---still tracks California's path without AB 1228 in 2024--2026,
when the general minimum rose more slowly.

The data would contradict the assumption if (i) the selected candidate's placebo RMSE is near or
above the 2 percent bound of hypothesis H2; (ii) the event path shows systematic gaps over
2022Q1--2023Q3; (iii) employment outside 722513 shows an effect; or (iv) the estimate moves
materially with the base or the donor pool. Check (ii) is an out-of-fit test only for plain DD: for
the trend-adjusted DD those seven quarters are the points the trend is fitted on, so the gaps there
test only curvature, and for synthetic control and SDID every non-COVID pre-period quarter is in the
fit window, so the path shows in-sample fit. For these candidates the out-of-fit pre-period evidence
is the selection rule's placebo RMSE. None of the four checks can confirm the assumption.

\subsection{Candidate estimators}
\label{sec:candidates}

Let $g_t = Y_{\mathrm{CA},t} - \bar Y_{\mathrm{donor},t}$ be the gap between California and a donor
average. The candidates are of four kinds.

\paragraph{DD.} $\hat\tau = \operatorname{mean}_{t\in P} g_t - g_B$.

\paragraph{Trend-adjusted DD (DDT).} In the style of \citet{clemens2026employment},
\begin{equation}
\hat\tau = \operatorname{mean}_{t\in P}\bigl[g_t - g_B - \hat b\,(t - B)\bigr],
\end{equation}
where $\hat b$ is the least-squares slope of $g_t$ on $t$ over the quarters $B-8, \dots, B$ outside
the COVID quarters: 2022Q1--2023Q3 (seven points) for the real base, nine points at the
pseudo-dates of the selection rule. This is equivalent to a state-specific linear trend in the
weighted two-by-two comparison.

For DD and DDT the donor average is formed in one of two ways: \emph{AGG}, the donor aggregate
series of summed employment, population and wages, logged, as Clemens and coauthors compare
California with the rest of the United States; or \emph{POPW}, the population-weighted mean of the
donors' log series with weights from population at $B$, as in \citet{dube2026labor}.

\paragraph{Demeaned synthetic control (SC).} Simplex weights $w$ minimise
\begin{equation}
\sum_{t\in F}\Bigl[\bigl(Y_{\mathrm{CA},t} - \bar Y_{\mathrm{CA},F}\bigr)
  - \sum_j w_j\bigl(Y_{j,t} - \bar Y_{j,F}\bigr)\Bigr]^2,
\end{equation}
and $\hat\tau$ is the mean of the same demeaned gap over $P$ \citep{doudchenko2016balancing,ferman2021synthetic}.
The demeaned form is needed because California's level of log employment lies outside the donors'
convex hull. With no penalty, the weights---and so the post-period estimate---are not unique when
the fit window is shorter than the number of donors, which happens at every pseudo-date and for R50
in the real window; the estimator is then the deterministic solve from uniform weights.

\paragraph{Synthetic DID (SDID).} We implement Algorithm~1 of \citet{arkhangelsky2021synthetic}
without sparsification. Unit weights $(\omega_0, \omega)$, with $\omega$ in the simplex, minimise
$\sum_{t\in F}\bigl(\omega_0 + \sum_j \omega_j Y_{j,t} - Y_{\mathrm{CA},t}\bigr)^2 + \zeta^2|F|\,\lVert\omega\rVert^2$,
with $\zeta = (1\cdot 8)^{1/4}\hat\sigma$ and $\hat\sigma^2$ the variance of the donors' first
differences over pairs of consecutive quarters inside $F$, as in their equation~(2.2). Time weights
$(\lambda_0, \lambda)$, with $\lambda$ in the simplex, minimise
$\sum_j\bigl(\lambda_0 + \sum_{s\in F}\lambda_s Y_{j,s} - \operatorname{mean}_{t\in W} Y_{j,t}\bigr)^2$
plus the negligible penalty $(10^{-6}\hat\sigma)^2 J\lVert\lambda\rVert^2$, for the estimation window
$W$. Then
\begin{equation}
\hat\tau = \Bigl(\operatorname{mean}_{t\in W} Y_{\mathrm{CA},t} - \sum_{s}\lambda_s Y_{\mathrm{CA},s}\Bigr)
  - \sum_j\omega_j\Bigl(\operatorname{mean}_{t\in W} Y_{j,t} - \sum_s\lambda_s Y_{j,s}\Bigr).
\end{equation}
$\omega$ and $\zeta$ are fitted once with the eight-quarter post window and reused for the yearly and
anticipation windows; $\lambda$ is refitted for each window.

All weights are exact minimisers of these objectives. On the Proposition 99 data our DD and SDID
estimates round to the values \citet[Table~1]{arkhangelsky2021synthetic} print ($-27.3$ and
$-15.6$), while our exact synthetic control differs from the printed $-19.6$ by about 0.1 packs per
capita, a difference that the reference package's own iterative solver and its default
sparsification reproduce; our design fixes the exact minimiser. For the wage outcome each estimator
is re-applied to log weekly wages with the same pool, refitting the weights (SC, SDID) or the trend
(DDT) on wages.

\paragraph{The candidate set.} Crossing DD and DDT with AGG and POPW weighting, the outcomes $E$
and $P$, and the pools D20 and R50 gives 16 candidates; SC and SDID, each with $E$ and $P$ and D20
and R50, give eight more. The 24 candidates, identified as \cand{DDT-POPW-P-D20} and so on, are
ordered DD, DDT (each AGG then POPW, $E$ then $P$, D20 then R50), SC, SDID, and this order is the
last tie-break. The set covers both published approaches: the detrended aggregate comparison with
the rest of the United States of Clemens and coauthors, and the per-capita SDID with federal-minimum
donors of Dube.

\subsection{The selection rule}
\label{sec:select}

The primary specification is the candidate that best reproduces California's own path at placebo
dates before COVID. The run that applies the rule reads only the panel truncated at 2019Q4 and is
the first analysis run after the design gate; it is run once. For each candidate $c$ and each
pseudo-treatment date $B' \in \{\text{2016Q1}, \dots, \text{2017Q2}\}$ (six dates), California is
treated at $B'$, the candidate's pool supplies the donors, and the windows $F(B')$, $A(B')$ and
$P(B')$ are defined as above, so every pseudo post window ends by 2019Q4 and contains no COVID
quarter. The first date is 2016Q1 because DDT needs $B'-8 \geq$ 2014Q1, the last 2017Q2 because
$B'+10 \leq$ 2019Q4. The pseudo-ATT $\hat\tau_c(B')$ is read in log points of employment, also for
per-capita candidates, and
\begin{equation}
\mathrm{RMSE}_c = \Bigl(\tfrac{1}{6}\textstyle\sum_{B'}\hat\tau_c(B')^2\Bigr)^{1/2}.
\end{equation}
The primary is the candidate with the smallest $\mathrm{RMSE}_c$. If RMSEs are equal after rounding
to four decimals, the smaller maximum $|\hat\tau_c(B')|$ wins, then the earlier candidate in the
order above. The rule's parameters were fixed on November 21, 2026 without computing any placebo
or any statistic of the outcomes, and no alternative set of pseudo-dates is computed. The paper
reports all 24 RMSEs.

Three properties of the rule should be kept in mind. First, $\hat\tau_c(B')$ is a fit error relative
to California's path without AB 1228, not an estimate of a zero policy effect: California's
2016--2019 path contains its general minimum-wage ramp, which the D20 donors lack entirely and the
R50 donors partly, so the rule favours candidates that track that ramp, such as DDT or R50
candidates. An R50 or DDT winner may reflect that shared ramp. Second, the rule has little
information: the six pseudo post windows overlap by three to seven of their eight quarters, so they
amount to about two independent draws, in calm years, with fit windows of 9 to 14 quarters against
31 in the real analysis, and the trend windows of DD and DDT are alike at the pseudo-dates but differ
after COVID. The winner may not be the best estimator for 2024--2026. Third, for SC candidates the
non-unique weights make the pseudo-ATTs depend on the solver; the run reports, without using it,
the largest gap between each SC candidate's pseudo-ATTs and those of an independent solver. For
these reasons the two published specifications are always reported as named secondary analyses and
can qualify the headline (Section~\ref{sec:decision}).

\subsection{Inference with one treated state}
\label{sec:inference}

\paragraph{Placebo estimates and the standard error.} For a statistic $\hat\tau_u$ of a treated
unit $u$---California in every post-period run---we compute the placebo estimate $\hat\tau_j$ of
each of the $J$ donors in $u$'s pool, treating $j$ as treated with the other $J-1$ donors as its
donors ($u$ excluded) and the same estimator, windows and outcome. Two standard errors are formed:
\begin{equation}
\SEs = \Bigl(\frac{1}{J}\sum_j \hat\tau_j^2\,\frac{\mathrm{pop}_j}{\mathrm{pop}_u}\Bigr)^{1/2},
\qquad
\SEn = \Bigl(\frac{1}{J}\sum_j \hat\tau_j^2\Bigr)^{1/2},
\qquad
\mathrm{SE} = \max(\SEs, \SEn),
\end{equation}
with populations at the base quarter. $\SEs$ is the population scaling of
\citet{ferman2019inference} as implemented by \citet{dube2026labor}. The maximum is conservative for
California. Suppose $\operatorname{var}(\hat\tau_j) = A + B/\mathrm{pop}_j$ with $A, B \geq 0$, a
state-level component and a size-related component. Then
$E[\SEs^2] = A\,\overline{\mathrm{pop}}/\mathrm{pop}_u + B/\mathrm{pop}_u$ and
$E[\SEn^2] = A + B\,\overline{1/\mathrm{pop}}$, against the truth $A + B/\mathrm{pop}_u$. Because
California is larger than every donor, $\SEs$ is at or below the truth and $\SEn$ at or above it for
every $A, B \geq 0$; scaling alone imposes $A = 0$, which is unlikely for a near-census like the QCEW.
The argument fails if the variance rises with size ($B < 0$); the level check below looks for that.
Since every ratio $\mathrm{pop}_j/\mathrm{pop}_{\mathrm{CA}}$ is below one, $\SEs < \SEn$ in every run,
and California's standard error is $\SEn$. As a diagnostic that decides nothing, we also report
Ferman and Pinto's own step: the least-squares regression of $\hat\tau_j^2$ on a constant and
$1/\mathrm{pop}_j$, giving $\hat A$, $\hat B$ and
$\mathrm{SE}_{\mathrm{FP}} = \max(0, \hat A + \hat B/\mathrm{pop}_u)^{1/2}$. Standard errors clustered
by state are not used.

\paragraph{Degrees of freedom.} California's $\SEn$ is the root mean square of $J$ placebo
estimates whose variances $s_j^2$ differ across donors, so its effective degrees of freedom are
fewer than $J$: the Satterthwaite value $(\sum_j s_j^2)^2/\sum_j s_j^4$ equals $J$ only when the
$s_j^2$ are equal. Every label uses Student's $t$ at $\nu = \min(J, \hat\nu)$, never at $J$, where
$\hat\nu$ is that formula with $s_j^2$ the mean over the six pseudo-dates of donor $j$'s squared
placebo estimate of the same statistic in the level check, computed on data ending in 2019Q4. So
$\nu$ is fixed for every statistic before any post-2023Q3 outcome is read. With about two
independent draws per donor the estimated $s_j^2$ are dispersed, which biases $\hat\nu$ downward
(to roughly $J/2$ when the true variances are equal), and when donor variances differ because of
size, $\SEn$ is already conservative for California and the small $\hat\nu$ widens the interval a
second time. Both costs are accepted, not corrected. So that readers can see how much of the
interval width is due to this penalty, the minimum detectable effect and the equivalence half-width
of H2 are also reported with $t_J$; no label uses those values.

\paragraph{Level check and the decision set.} A run on the truncated panel, after the design gate
and before any run that reads a post-2023Q3 outcome, measures how the procedure behaves where the
effect is known to be absent. At each of the six pseudo-dates, every unit $u$ in the pool and
California in turn is treated. California uses the whole pool as donors; every other unit uses the
pool without itself, with California excluded from both its donors and its placebos, so that no
placebo-unit test carries California's 2014--2019 minimum-wage ramp. For each statistic the run
reports rejection rates of $T_u = \hat\tau_u/\mathrm{SE}_u$ at the nominal 5 percent one- and
two-sided levels, with $\mathrm{SE}$, $\SEs$ and $\SEn$ separately: over all tests, by population
tercile, over the decision set $L$, and for California's six tests. The statistics checked are the
selected candidate's pooled wage, pooled employment and year-two-minus-year-one statistics, the
pooled employment statistics of the secondary analyses S1 and S2, and the selected candidate's
pooled employment statistic with the other pool.

The decision set $L$ is the ten most populous D20 states at each pseudo-date (60 tests). Their
binding minimum was the federal \$7.25 throughout 2014--2019, so their pseudo-ATTs have no
own-policy content, unlike California's and many R50 states'. Because the six pseudo post windows
overlap, $L$ holds about 20 independent draws (all tests: about 42 for D20 and 102 for R50). $L$
contains no unit as large as California, so it can reveal over-rejection that grows with size but
cannot measure California's own level; for California the level is argued from the conservative
standard error and critical value, not measured, and we say so wherever we report it.

\paragraph{Critical values and intervals.} For each statistic,
\begin{align}
c_{\mathrm{hi}}(p) &= \max\bigl(t_{\nu,p},\; Q^{\mathrm{higher}}_p(T_L)\bigr), \quad p \in \{0.95, 0.975\},\\
c_{\mathrm{lo}}(p) &= \min\bigl(-t_{\nu,1-p},\; Q^{\mathrm{lower}}_p(T_L)\bigr), \quad p \in \{0.025, 0.05\},
\end{align}
where $Q_p(T_L)$ is the empirical $p$-quantile of the 60 values of $T$ in $L$ (taken as the next
higher or lower observation). With 60 values the 0.025 and 0.975 quantiles are the second smallest
and second largest, attainable but noisy; that noise can only widen an interval relative to
$t_\nu$. The intervals are
\begin{equation}
\text{95\%: } \bigl[\hat\tau - c_{\mathrm{hi}}(0.975)\,\mathrm{SE},\; \hat\tau - c_{\mathrm{lo}}(0.025)\,\mathrm{SE}\bigr],
\qquad
\text{90\%: } \bigl[\hat\tau - c_{\mathrm{hi}}(0.95)\,\mathrm{SE},\; \hat\tau - c_{\mathrm{lo}}(0.05)\,\mathrm{SE}\bigr],
\end{equation}
and the one-sided 5 percent test of H1 rejects if $\hat\tau/\mathrm{SE} > c_{\mathrm{hi}}(0.95)$.
$p$-values from $t_\nu$ are reported for information only; every label uses the intervals. With
$J = 20$ or $50$ donors a 5 percent test is attainable without California ranking first among the
placebos. The selected candidate's critical values serve H1, H2, H3 and the anticipation estimate;
S1 and S2 have their own, as does the other-pool robustness check; the remaining robustness and
placebo runs and the event-path bands use the selected candidate's employment critical values with
$t$ at each run's own $\nu$.

\paragraph{The SDID floor.} Before the design gate, a Monte Carlo study checked the inference on
synthetic panels with the real 2016 populations of California and the donors and a true effect of
zero, with state-level shocks, size-related noise, or both, with and without autocorrelation, and
in two designs where the state-level shock variances differ: five oil states with four times the
variance and California at the donors' mean, and the same with California at twice that mean. For
DD, DDT and synthetic control, California's tests at $t_J$ with the maximum standard error kept
their level in the designs with equal shock variances. SDID's placebo $t$ statistic did not: in one
cell (D20 donors, pure state-level shocks, no autocorrelation) its two-sided rejection rate at
$t_J$ exceeded the pre-set bound, and no error in the code was found; the excess lies in the tails
of SDID's placebo $t$, not in the scale of its standard error. Every SDID statistic therefore has a
calibrated floor on its critical values: $c_{\mathrm{hi}}(p)$ is also at least the
$p$-quantile, and $c_{\mathrm{lo}}(p)$ at most the $p$-quantile, of California's $t$ statistic in
fresh pure-state-shock replications for the pool the run uses, with and without autocorrelation,
the more extreme of the two. The floor applies to S2 and, if the rule selects an SDID candidate, to
every statistic of the primary. It comes from the geometry of the pseudo-dates, not from the real
fit window, so it transfers to the real runs by assumption; it can only widen an SDID interval. In
the design with unequal shock variances one SDID cell (D20, no autocorrelation) rejected above its
bound even at the Monte Carlo degrees of freedom, again with no error found; we report that rate as
the size of an SDID test with D20 under that design. The design with California twice as noisy
as its donors shows how far California's rate can exceed 5 percent if its own noise is larger than
its size implies; it cannot tell whether it is. The rates and the floor values are reported with
the results.
% The rates (D20 pure-A uncorrected and with the floor; hetero-A and hetero-A-CA2 for SDID) and the
% q_cal values enter as macros of claims backed by level-mde (claims_plan item 1).

\subsection{Power}
\label{sec:power}

The same truncated-data run computes the power of the design. For each statistic, $\mathrm{SE}_{\mathrm{pre}}$
is the root mean square of California's standard error over the six pseudo-dates. The minimum
detectable effects at 80 percent power are
\begin{equation}
\mathrm{MDE}_{\mathrm{H2}} = \bigl(c_{\mathrm{hi}}(0.975) + t_{\nu,0.80}\bigr)\mathrm{SE}_{\mathrm{pre}},
\qquad
\mathrm{MDE}_{\mathrm{H1}} = \bigl(c_{\mathrm{hi}}(0.95) + t_{\nu,0.80}\bigr)\mathrm{SE}_{\mathrm{pre}},
\end{equation}
with $\mathrm{MDE}_{\mathrm{H3}}$ defined like $\mathrm{MDE}_{\mathrm{H2}}$. For the size label of H2,
$h_{90} = \max\bigl(c_{\mathrm{hi}}(0.95), -c_{\mathrm{lo}}(0.05)\bigr)\mathrm{SE}_{\mathrm{pre}}$ is the
half-width of the 90 percent interval; the label ``no effect as large as 2\%'' is reachable only if
$h_{90} < 0.02$. The run also computes the probability of each of the nine H2 labels when
$\tau = 0$ and when $\tau = -0.035$, close to the two-year median of CEMN, from
$\hat\tau = \tau + \mathrm{SE}_{\mathrm{pre}}T$ with $T \sim t_\nu$. Because $\mathrm{SE}_{\mathrm{pre}}$
rests on 2014--2019 noise and on seasonal factors fitted to the same years, it is a lower bound; the
post-period runs recompute $\mathrm{MDE}_{\mathrm{H2}}$, $h_{90}$ and $\mathrm{MDE}_{\mathrm{H3}}$ with
the real standard error, for information, and the rule for H3 uses the larger of the two
$\mathrm{MDE}_{\mathrm{H3}}$ values.

\subsection{Hypotheses and decision rules}
\label{sec:decision}

\paragraph{H1, first stage in weekly earnings.} The pooled 2024Q2--2026Q1 effect on the log average
weekly wage is positive: one-sided test at 5 percent, also reported for each year. Weekly earnings
per job are not the hourly wage the law sets; fewer hours would bias the effect toward zero. If H1
is not supported, H2 is reported as uninformative about the wage floor.

\paragraph{H2, employment.} With $b = 0.02$ log points, let $[L, U]$ be the 95 percent and
$[l, u]$ the 90 percent interval of the pooled employment effect. The sign is negative if $U < 0$,
positive if $L > 0$, and undetermined otherwise. The size is smaller than $b$ if $-b < l$ and
$u < b$ (two one-sided tests at 5 percent), at least $b$ if $u < -b$ or $l > b$, and undetermined
otherwise. The headline label is the pair, from the exhaustive and mutually exclusive
Table~\ref{tab:labels}.

\begin{table}[tbp]
\centering
\caption{Labels of hypothesis H2}
\label{tab:labels}
\small
\begin{tabularx}{\linewidth}{@{}l>{\raggedright\arraybackslash}X>{\raggedright\arraybackslash}X>{\raggedright\arraybackslash}X@{}}
\toprule
& \multicolumn{3}{c}{Size (90\% interval against $b = 0.02$)} \\
\cmidrule(l){2-4}
Sign (95\% interval) & at least $b$ & smaller than $b$ & undetermined \\
\midrule
negative & decline of at least 2\% & decline smaller than 2\% & decline, size not distinguishable from 2\% \\
\addlinespace
undetermined & inconclusive & no effect as large as 2\% & inconclusive \\
\addlinespace
positive & increase of at least 2\% & increase smaller than 2\% & increase, size not distinguishable from 2\% \\
\bottomrule
\end{tabularx}
\end{table}

The bound was chosen to separate the two published readings; Section~\ref{sec:res-power} shows
that at this design's precision it cannot. The published pooled weekly-wage effects are about
6 percent \citep{clemens2026employment} to 7 percent \citep[Table~1, Panel~A]{dube2026labor}, close
to but not the same estimand as ours; at those wage effects an employment effect of $-0.02$ is an
own-wage elasticity of about $-0.29$ to $-0.33$. The two-year QCEW estimates of CEMN range from
$-2.02$ to $-5.03$ percent, so their least negative sits at the bound: an interval inside $\pm0.02$
excludes their median and detrended estimates but not every one of their specifications, and an
interval wholly below $-0.02$ excludes Dube's SDID estimate. A ``decline of at least 2\%'' would
therefore count against the near-zero reading; ``no effect as large as 2\%'' or ``decline smaller
than 2\%'' would count against the reading of a decline of about 3.5 percent.

\paragraph{H3, year two against year one.} The difference $\tau_2 - \tau_1$, with the same estimator
and inference, is labelled ``larger decline in year two'' if its 95 percent interval is below zero,
``smaller decline in year two'' if above, and ``no detectable change'' otherwise. If
$\mathrm{MDE}_{\mathrm{H3}}$ exceeds 0.04, twice $b$, H3 is reported as uninformative whatever the
estimate. The difference is read next to the yearly wage effects and the implied yearly
elasticities, because the bite of the law fell in year two and a change cannot be attributed to
time alone.

\paragraph{Elasticities.} The own-wage elasticity, pooled and by year, is the employment effect
divided by a wage effect computed with the employment fit reused (for SC the employment weights, for
SDID the employment unit and time weights; DD and DDT have no fitted weights). Numerator and
denominator then use one comparison, so dilution by uncovered establishments roughly cancels. For
SC and SDID the ratio with the refitted wage effect of H1 is reported next to it. An elasticity is
computed only for windows whose H1 wage effect rejects and whose reused-weights wage effect is
positive, and only as a point value.

\paragraph{Qualifications of the headline.} A final run applies these rules mechanically. Each
secondary analysis and robustness check also reports the label the H2 rule would give, marked as
not the headline. If S1, S2, the 2023Q4 base, the pre-AB 257 base, the other pool or the
final-data window gives a different sign category from the primary, the headline is stated as not
robust to that check. If the 95 percent interval of the placebo outcome (employment outside 722513)
excludes zero, ``not separable from a California-wide shock'' is appended to the H2 label. The
abstract and the first results paragraph report the primary, S1 and S2 labels side by side.

\paragraph{Multiple testing.} No family-wise correction is applied: H2 is the single primary
employment test, H1 a gate for reading it, and H3 a separately labelled secondary with its own power
rule. Secondary analyses, robustness checks and placebos never create a new finding. With two
placebo outcomes, the chance that one of them ``fails'' at 95 percent by chance alone is about 0.10.

\subsection{Secondary analyses}
\label{sec:secondary}

Two named secondary analyses re-estimate each published approach with our timing and inference,
whatever the rule selects.

\paragraph{S1, the detrended comparison of Clemens and coauthors.} With
$g_t$ the gap between California's seasonally adjusted log employment and that of the R50
aggregate,
\begin{equation}
\hat\tau_{\mathrm{S1}} = \operatorname{mean}_{t\in P}\Bigl[g_t - g_B - (t-B)\,\frac{g_B - g_{B-8}}{8}\Bigr],
\end{equation}
which keeps their endpoint 2021Q3--2023Q3 growth trend; its start, 2021Q3, is the only COVID
quarter used in any estimate. The undetrended version is reported too. Inference uses R50 and S1's
own critical values, alongside their placebo rank $p$-value $(b+1)/(m+1)$. As a cross-check against
print, the run also computes their September-to-September fixed-horizon estimate for 2023--2025
on third-month levels with their 48-state pool, simple and detrended, next to
\citet[Table~3, Panel~B]{clemens2026employment} ($-2.74$ and $-3.72$).

\paragraph{S2, the SDID of Dube.} SDID on seasonally adjusted log employment per capita with the D20
pool, a fit window of 2015Q1--2023Q3 without the COVID quarters and the post window 2024Q2--2026Q1;
the wage effect reuses the employment weights, as Dube does. Inference uses S2's own critical
values, which include the D20 SDID floor. As a cross-check, the run also applies Dube's own timing
(fit window to 2023Q4, 2024Q1 dropped, post window 2024Q2--2025Q3), next to
\citet[Table~1, Panel~A]{dube2026labor} ($-0.26$ with standard error 0.54; wage effect 6.99).
Because our SDID follows Algorithm~1 with an intercept and separate penalties and uses our own
seasonal factors, and Dube's code is not available to us, that cross-check measures the combined
difference.

\subsection{Robustness and placebo checks}
\label{sec:robust-plan}

Each of the following uses the selected estimator, pool and outcome unless stated, and is reported
whatever it shows.
\begin{enumerate}
  \item \cand{rob-base-2023q4}: Dube's base, $B$ = 2023Q4, with 2024Q1 dropped and the same post
        window.
  \item \cand{rob-pre-ab257}: $B$ = 2022Q2, the last quarter before AB 257 was enacted, with
        2022Q3--2024Q1 excluded from the base, the fit and the post window. For DDT the trend uses
        the last nine non-COVID quarters up to 2022Q2. This reports the base both ways and also
        covers anticipation of AB 1228 from mid-2023, which \citet{sosinskiy2026effects} discuss.
  \item \cand{rob-pool}: the other donor pool (D20 for R50 and the reverse).
  \item \cand{rob-7225}: all restaurants (7225) as outcome. The four donors with a suppressed 2024Q4
        row are dropped and listed, leaving 18 D20 or 46 R50 donors, with $t$ at the run's own
        $\nu$; a second number keeps all donors and drops 2024Q4 from every post window. A 722513
        decline with no 7225 decline would point to reclassification between restaurant codes or
        substitution toward other restaurants rather than job loss in restaurants.
  \item \cand{rob-final-quarters}: the post window 2024Q2--2025Q4, without the preliminary 2026Q1,
        to show how much the one new quarter moves the estimate.
  \item \cand{plc-nonff} (placebo outcome): employment in all covered industries except 722513,
        of which the law's direct share is small. A 95 percent interval that excludes zero means the
        comparison does not hold for California's aggregate labour market, and qualifies the H2
        label; no re-specification follows.
  \item \cand{plc-722511}: full-service employment, reported only. It is not a clean placebo:
        covered chains are also coded in 722511 \citep{melo2026covered}, and wage and price
        spillovers to full service are reported \citep{clemens2026prices}. A decline of similar size
        to H2 would mean the result is not specific to the limited-service code; a null or
        opposite-signed effect does not validate the design, and 722511 is never used as a
        triple-difference comparison.
  \item \cand{anticip}: the employment and wage effects over 2023Q4--2024Q1 against the 2023Q3
        base, without a label or test.
  \item \cand{event-path}: the quarterly gap between California and the selected counterfactual,
        2014Q1--2026Q1, for employment and wages, with pointwise bands; COVID quarters shaded, AB
        257 (2022Q3) and AB 1228 (2023Q4) marked. The 2022Q1--2023Q3 gaps are an out-of-fit check
        only if plain DD is selected; for DDT they are residuals around the fitted trend, and for SC
        and SDID the whole pre-period is in-sample fit. The figure is not a pre-trend test for those
        candidates, and its note says which case applies.
\end{enumerate}

\subsection{Pre-specification and what was known}
\label{sec:prespec}

The runs follow a fixed order: the selection rule first, on the truncated panel; then the level
check and power run, also on the truncated panel; then every run that reads an outcome after 2023Q3.
Each post-period run reads the selected candidate and the critical values from the outputs of the
first two runs; nothing is typed in by hand. If a reviewed bug fix forces the selection run to be
repeated, the superseded record is kept and the paper reports every selection run and its winner.

The study is pre-specified but not blind. The design's choices were made knowing the published map
of how each specification choice moves the estimate (Section~\ref{sec:lit-map}), which is why they
are delegated to a rule and both published specifications are reported. The PI and the literature
scout also read the public descriptive series of \citet{steier2026recordinglaw}, including its
March 2026 employment index values for California and the rest of the United States and its 2026Q1
weekly-wage index values (Section~\ref{sec:lit-estimates}), before the design was fixed. The roles
that analyse the data computed nothing on outcomes after 2023Q3 before the design gate. The design
was reviewed in \CrewRoundsDesign{} rounds by a reviewer agent that did not author it. The rules
that settle the outcomes of the code check (among them the SDID floor and the Proposition 99
synthetic-control check), the adoption of the implementation choices, and two report-only additions
were made after the last of those rounds; they were not
reviewed by a design round, and the human PI approved the design with them on record.

\subsection{Threats to validity}
\label{sec:threats}

\paragraph{Search moved into the rule.} Fixing a selection rule moves the specification search into
the rule's parameters: pseudo-dates, windows, COVID handling, outcome scale, tie-break, pools,
inference and cutoff. All were fixed before any placebo was computed, the selection runs once on
data ending in 2019Q4, and both camps' specifications are reported. The rule itself has little
information and may favour candidates that track California's own minimum-wage ramp
(Section~\ref{sec:select}); there is no re-selection.

\paragraph{COVID, AB 257 and anticipation.} Every eight-quarter pseudo post window starting from
2018Q3 would contain COVID quarters, so all pseudo-windows end by 2019Q4; the real fit windows
straddle the COVID gap, which the pseudo-dates do not mimic. The quarters 2022Q3--2023Q3, including
the base and the end of every real fit window, fall under the pending AB 257
(Section~\ref{sec:ab257}); the selection, level check and power calculations avoid them, and the
pre-AB 257 base reports the base both ways. Anticipation of AB 1228 before signing, which
\citet{sosinskiy2026effects} discuss from as early as July--August 2023, would bias a 2023Q3 base; the
pre-AB 257 base and the event path speak to it. Anticipation between signing and the effective date
is estimated separately and checked with Dube's base.

\paragraph{Inference with one treated state.} The maximum standard error is conservative for the
largest unit only if variance does not rise with size; California-specific noise beyond its size
would make the tests over-reject, and the Monte Carlo design with California twice as noisy shows by
how much, without saying whether it is so. $L$ contains no unit as large as California, so its level
for California is argued, not measured. The degrees-of-freedom rule and the SDID floor widen the
intervals, and the power quantities show the cost. The minimum detectable effect from 2014--2019
noise, with seasonal factors fitted to the same years, may understate post-COVID noise; the real
standard error is reported next to it. Stale seasonal factors after COVID mostly average out over
two full years, except at the single third-quarter base, which the 2023Q4 base checks.

\paragraph{Comparison states and concurrent shocks.} R50 states that raised their own minimum wages
in 2024--2026 bias an R50 estimate toward a smaller decline; D20 states had no increase but differ
economically. The other-pool check and S1 against S2 show the difference. In California, the general
minimum wage rose, affecting uncovered 722513 and 722511 workers; a health-care minimum wage took
effect from October 2024 \citep{clemens2026employment}; and population grew more slowly than in the
rest of the country \citep{clemens2026employment}, which matters for per-capita candidates and for the extrapolated population of
2025Q3--2026Q1. The placebo outcome is the check, and its failure qualifies H2.

\paragraph{Measurement.} The treated sector is wider than and different from 722513
(Section~\ref{sec:coverage}): covered chains sit also in 722515, never used, and 722511, a non-clean
placebo only; classification may respond to the policy, which the 7225 check addresses. The estimand
is diluted by uncovered establishments, and no rescaling is done; the elasticity is roughly invariant
to dilution. If QCEW 722513 tracks covered-chain employment poorly across states and over time, as
\citet{melo2026covered} argue, the counterfactual can be distorted in ways the placebo standard error,
which samples only donor divergence present in 2014--2019, does not capture; the equivalence verdict
of H2 is a statement about QCEW 722513. The QCEW counts jobs in the pay period including the 12th,
not hours, so lower turnover can lower counts \citep{dube2026labor}; weekly earnings are not hourly
wages, and nominal wages carry differential price changes. The 2026Q1 data are preliminary; the
final-data window drops them.

\paragraph{Code errors.} The estimator code was checked on the Proposition 99 data against
\citet[Table~1]{arkhangelsky2021synthetic}, on synthetic identities and in the Monte Carlo study
before any analysis run; it is reviewed before it runs, cross-checked against the published
estimates of both camps, and rerun from scratch in a clean environment before submission.

\subsection{Departures from the track rules and the published specifications}
\label{sec:departures}

\begin{enumerate}
  \item Standard errors are not clustered at the level at which treatment varies: with one treated
        state, cluster-robust standard errors are invalid \citep{dube2026labor,clemens2025did,sosinskiy2026effects}.
        Placebo-based inference with a level check fixed in advance is used instead.
  \item S1 adapts the monthly fixed-horizon estimator of Clemens and coauthors to a pooled quarterly
        average, which needs the seasonal adjustment their September-to-September comparison does
        not; the faithful fixed-horizon number is reported as a cross-check. S1 also keeps their
        endpoint trend, whose start is a COVID quarter.
  \item S2 uses our implementation of Algorithm~1 of \citet{arkhangelsky2021synthetic}, with an
        intercept and separate penalties, and our 2014--2019 seasonal factors, not Dube's code;
        \citet{clemens2026synthetic} report that Dube's SDID omits the intercept and uses one penalty.
  \item R50 keeps Alaska and Michigan, which CEMN drop; the cross-check against their print uses
        their 48-state pool.
\end{enumerate}
